La aplicación de métodos cuánticos a problemas de optimización financiera representa una línea de investigación activa en computación cuántica. Aunque la optimización de portafolios ha sido extensamente abordada mediante técnicas clásicas, la posibilidad de codificar información financiera en espacios de mayor dimensionalidad que habilita la computación cuántica ofrece un enfoque alternativo que merece evaluación experimental rigurosa.

Este trabajo propone la implementación y validación experimental del algoritmo Quantum Hierarchical Risk Parity (QHRP), que integra codificación cuántica con métodos de ordenamiento jerárquico de riesgos. El objetivo es evaluar si las representaciones cuánticas del espacio de correlaciones entre activos pueden reproducirse y construirse de forma consistente, permitiendo una comparación metodológica con el enfoque clásico HRP bajo condiciones controladas.

En este documento se detalla la formulación teórica del circuito cuántico empleado como ansatz y el protocolo experimental para la comparación con métodos clásicos. Se enfatiza la reproducibilidad de los cálculos, la consistencia de la codificación de datos y las limitaciones observadas en la validación numérica. Los hallazgos contribuyen a una mejor comprensión de la viabilidad y las restricciones de estos métodos en aplicaciones financieras prácticas.

\hfill \begin{flushright}Iker González\end{flushright}\hfill
